\section{Introduction}
\label{sec:intro}

\begin{figure}
    \centering
    \includegraphics[width=1.0\linewidth]{Images/figure_1.pdf}
    \caption{\textbf{KeySync's contributions.} Unlike existing methods, KeySync generates high-resolution lip-synced videos that are closely aligned with the driving audio while minimizing leakage from the input video and seamlessly handling facial occlusions.}
    \label{fig:intro_figure}
\end{figure}

Audio-driven facial animation has recently seen substantial progress with the introduction of new generative models such as Generative Adversarial Networks (GANs)~\cite{goodfellow2020generative, Vougioukas, zhou2019talking} and diffusion models~\cite{ho2020ddpm, stypulkowski2023, xu2024hallohierarchicalaudiodrivenvisual, wang2024V-Express, chen2024echomimiclifelikeaudiodrivenportrait}. In contrast, the adjacent field of lip synchronization (also known as lip-sync) has experienced comparatively slower advancements~\cite{guan2023stylesync, DINet, Prajwal}. This disparity is surprising given that lip-sync has similar applications, ranging from facilitating multilingual content production to enhancing virtual avatars~\cite{zhen2023human, zhan2023multimodal}. A potential reason for this slower progress is that while lip synchronization may seem like a simpler task than animating the full face from audio, it presents unique challenges that remain largely unaddressed.

One of the primary limitations of current methods is their low-resolution output, typically constrained to 256$\times$256, which has become a de facto standard. While this resolution may be computationally efficient, it significantly hinders real-world applicability, where higher resolution outputs are necessary for practical deployment. Another important limitation is that these methods struggle to maintain temporal consistency across generated frames. Most state-of-the-art approaches are frame-based, requiring additional mechanisms to enforce consistency, while other methods adopt a two-stage framework, where motion modelling is separated from pixel-level synthesis~\cite{maketalk, DiffDub, stylepres}. However, such decompositions do not inherently guarantee smooth transitions between frames, as each frame is generated independently, often leading to temporal discontinuities. 

Other approaches attempt to enforce temporal coherence using pre-trained perceptual models~\cite{latentsync} or an additional sequence discriminator~\cite{diff2lip}. Nevertheless, these methods offer only indirect control over frame-to-frame consistency, often resulting in subtle visual artifacts and unnatural mouth movements that degrade realism and limit practical usability. Finally, another proposed solution has been to condition generation on past frames~\cite{michaldub} to ensure that newly generated frames remain consistent with previous ones. However, this technique is prone to error accumulation over long sequences, further degrading temporal stability.

Beyond temporal consistency, a key but often overlooked issue is expression leakage, where models inadvertently retain facial expressions from the original input in the generated video. Regrettably, most existing works focus excessively on lip synchronization as a reconstruction task on paired audio-visual data, and neglect the cross-synchronization scenario, where a non-matching audio clip is used to re-animate the original video. As a consequence, they typically exhibit major expression leakage from the original video, severely degrading the synchronization between the generated video and the input audio in the latter scenario. Notably, this behaviour jeopardizes the viability of these models for applications such as automated dubbing, where the audio and video are naturally mismatched. 

To alleviate the issue of expression leakage, different masking strategies have been devised. Some methods mask only the mouth region while preserving facial areas such as the jaw and cheeks from the original videos, potentially leading to leakage since these regions also convey information about mouth movements~\cite{StyleLipSync, DINet}, while others adopt broader masks that risk discarding important contextual cues~\cite{zhang_musetalk_2024, VideoReTalking}. Remarkably, the impact of these masking strategies on generalization and robustness remains largely unexplored, and no consensus exists on the optimal approach. Lastly, another potential complication lies in occlusion handling. Most existing models assume an unobstructed view of the mouth, whereas, in the real world, occlusions caused by hands, objects, or motion blur are frequent. In practice, this means that the lack of explicit occlusion-handling mechanisms significantly limits the applicability of current models. 

To address these challenges, we propose KeySync, a two-stage lip synchronization framework that leverages recent advances in facial animation~\cite{keyface} to generate high-fidelity videos with lip movements that are temporally consistent and aligned with the input audio. To minimize leakage from the input video, we devise a masking strategy that adequately covers the lower face while retaining the necessary contextual regions. Furthermore, we augment this mask by excluding facial occlusions using a video segmentation model~\cite{sam2}, resulting in a method that can consistently handle occlusions without uncanny visual hallucinations. Our primary contributions, illustrated in Figure~\ref{fig:intro_figure}, can be summarized as:
\begin{itemize}
    \item \textbf{State-of-the-art lip synchronization}: KeySync achieves state-of-the-art lip synchronization performance at a resolution of $(512\times512)$,  surpassing the common $(256\times256)$ standard. It outperforms all competing methods in terms of quality and lip movement accuracy according to several objective metrics and a holistic user study. We observe particularly noticeable improvements in the cross-synchronization setting (where there is a mismatch between the input video and audio), enabling promising real-world applications such as automated dubbing.
    \item \textbf{A new strategy for occlusion handling}: We propose a new inference-time strategy for occlusion handling by excluding occluding objects from our mask automatically using a pre-trained video segmentation model. Through qualitative and quantitative analysis, we show that this method is consistently effective in handling occlusions.
    \item \textbf{A novel leakage metric}: To the best of our knowledge, we propose the first lip synchronization leakage metric (LipLeak), which computes the ratio of non-silent frames generated from a silent audio and a non-silent video, effectively measuring how often the lip movements from the input video leak into the generated video.
\end{itemize}
